%=============================================================================!
% 3.0  CHAPTER OPENING                                                        !
%=============================================================================!
Chapter~2 predicted motion by solving the equation of motion: the forces
fix the acceleration, from which we obtain the velocity and position.
Some questions can be answered without finding the whole trajectory.
A bead sliding into a dip, for example, may have enough speed to climb the
next rise; the two carts of \cref{sec:ne-bodies} reach definite speeds
when their spring falls away. Energy allows us to determine those speeds
by comparing the starting and finishing states.

We begin with the work done by a force and show from Newton's second law
that it equals the change in kinetic energy. For a force that depends
only on position along a line, the work defines a potential energy.
The sum of kinetic and potential energy is then conserved, and a graph
of the potential shows the accessible positions, turning points and
possible equilibria. Friction transfers energy out of the motion we are
following. In two or three dimensions, position dependence alone no longer
guarantees a potential: the work must also be independent of the path.
We develop the line integral, partial derivatives, gradient and curl as
they become necessary to test this condition. Applied to atoms, these
ideas give a practical check on an isolated simulation with conservative
forces: its total energy should remain constant within the numerical
error. Chapter~4 uses the shape of the potential to explain oscillations.
